% !PRINTED 531
\weberpage{531}

\gesperrt{Welche Werthe man den mehrwerthigen algebraischen Grössen, die bei der Auflösung auftreten, auch beilegen mag, die Ausdrücke (15) stellen also immer die Wurzeln unserer biquadratischen Gleichung in irgend einer Reihenfolge dar.}

In Bezug auf die übrigen Wege zur Auflösung der biquadratischen Gleichung können wir uns kürzer fassen.

Wenn wir zunächst nicht $\sqrt{D}$ adjungiren, sondern eine zu der Gruppe $P_1$ gehörige Function, so erhalten wir eine cubische Resolvente, die nicht Normalgleichung ist. Wir können für diese Function etwa $y^2$ wählen, was der cubischen Gleichung
\weq{16}{y^6 - 2Ay^4 + A^2y^2 - D = 0}
genügt. Besser noch nimmt man als Wurzeln der cubischen Resolvente
\weq{17}{z = y_1 - y_2, \quad z_1 = y_2 - y, \quad z_2 = y - y_1.}
Es gehört dann $z$ zur Gruppe $P_1$, wie schon aus der zweiten Gleichung (10) hervorgeht, oder auch direct eingesehen wird. Aus (8) und (9) erhält man für $z$ die cubische Gleichung
\weq{18}{z^3 - 3Az + B = 0.}
Man erhält ferner, wenn man
\weq{19}{u = \alpha\alpha_1 + \alpha_2\alpha_3}
setzt,
\[
z = a_2 - 3u
\]
und daraus die Resolvente für $u$:
\weq{20}{u^3 - a_2u^2 + (a_1a_3 - 4a_4)u + 4a_2a_4 - a_1^2a_4 - a_3^2 = 0.}

Wenn man aber von der Gleichung (18) oder (20) nicht eine, sondern alle Wurzeln adjungirt, so gelangt man wieder zu der Gruppe $Q_1$. Die Grössen $v_1, v_2, v_3$ sind nach (12) unmittelbar durch $z, z_1, z_2$ ausgedrückt, und man kann also die Formeln (15) zur Darstellung der $\alpha$ auch nach diesem Verfahren anwenden.

Man kann ferner zur Lösung der biquadratischen Gleichung dadurch gelangen, dass man zuerst eine zu der cyklischen Gruppe $P_2$ gehörige (eine cyklische) Function der Wurzeln adjungirt, wie z.\,B.
\weq{21}{w = \alpha\alpha_1^2 + \alpha_1\alpha_2^2 + \alpha_2\alpha_3^2 + \alpha_3\alpha^2.}
Diese Function ist sechswerthig und ist also die Wurzel einer Gleichung sechsten Grades. Diese Gleichung sechsten Grades ist aber, wie man leicht sieht, imprimitiv, und kann auf
%CONT>>
